\tetePremStssMeth
\vspace*{-36pt}
\titreActivite{Notation scientifique et unités}

\vspace*{-10pt}
\begin{objectifs}
  \item Revoir les règles de calcul associées aux puissances de 10,
  voir ou revoir la notation scientifique.
\end{objectifs}


%%%%
\titreSection{Le système international de mesure}

%%
\begin{doc}{Le système international}
  Pour comparer des grandeurs entre elles, il faut les exprimer avec les \important{mêmes unités de mesures.} % exemple centime et euros
  
  Pour pouvoir communiquer facilement d'un pays à un autre, le \important{système international (SI)} a été développé par la Conférence Générale des Poids et Mesures (CGPM) pendant la révolution française.
  Le système international est composé de \important{sept unités de bases.}

  En physique on est amené à décrire des \important{échelles} très variées, par exemple quand on mesure la taille d'une molécule ($\sim \qty{e-9}{\metre}$), d'une cellule ($\sim \qty{e-6}{\metre})$ ou d'un humain ($\sim \qty{1}{\metre}$).
  
  \begin{importants}
    Pour simplifier la manipulation des grandeurs éloignées de l'unité, chaque \important{puissance de \num{1000}} est associée à un \important{préfixe} dans le système international.
  \end{importants}

  \begin{center}
    \begin{tblr}{
      hlines, vlines, row{1} = {couleurSec-100}, colspec = {c c c l}
    }
      Puissance  & Préfixe & Symbole & Nombre décimal \\
      \num{e9}   & giga    & G       & \num{1 000 000 000} \\
      \num{e6}   & mega    & M       & \num{1 000 000} \\
      \num{e3}   & kilo    & k       & \num{1 000} \\
      $10^0$     &         &         & \num{1} \\
      \num{e-3}  & milli   & m       & \num{0,001} \\
      \num{e-6}  & micro   & $\mu$   & \num{0,000 001} \\
      \num{e-9}  & nano    & n       & \num{0,000 000 001} \\
    \end{tblr}
  \end{center}

  C'est pour ça qu'on va avoir besoin de manipuler des puissances de dix et des unités de mesure.
\end{doc}


%%%%
\titreSection{Rappels sur les puissance de 10}

%%
\begin{doc}{Les puissances de 10}
  Les puissances indiquent qu'on va répéter une multiplication ($2^3 = 2 \times 2 \times 2 = 8$), par convention n'importe quel nombre puissance 0 est égal à 1 ($2^0 = 10^0 = 2798^0 = 1797431^0 = 3^0 = 1$). 

  Si la puissance est négative, on va répéter une division ($2^{-3} = 1 / 2^3 = 1 / (2\times 2\times 2) = 1/8$).

  Dans le cas des puissances de 10, on a donc plusieurs multiplications ou divisions successives par 10, ce qui permet de passer des unités, aux dizaines,  centaines, milliers, etc.

  \begin{importants}
    Concrètement, \important{chaque multiplication par 10 rajoute un zéro à un nombre et chaque division par 10 enlève un zéro.}
    Pour la division, si on descend en dessous de l'unité ``enlever'' des zéro consiste à rajouter des zéro après la virgule.

    \exemple*{} $10^{-2} = 0,01$ ; $10^{-4}= 0,0001$ ; $10^4 = \num{10000}$.
  \end{importants}
\end{doc}

\numeroQuestion Écrire les nombres correspondant aux puissances de 10 suivantes : \\
$\num{e2}  =$ \texteTrou[0.15]{\num{100}} \qq{}
$\num{e5}  =$ \texteTrou[0.15]{\num{100000}} \qq{}
$\num{e-3} =$ \texteTrou[0.15]{\num{0,001}} \qq{}
$\num{e-1} =$ \texteTrou[0.15]{\num{0,1}}

%%
%\pasCorrection{\vspace*{-18pt}}
\begin{doc}{Règles de calculs}
  Il y a deux règles de calculs à connaître pour les calculs de puissances
  \begin{importants}
    \pointCyan $10^a \times 10^b = 10^{a + b}$ 
    \qq{} \exemple $10^4 \times 10^5 = 10^{4 + 5} = 10^{9}$ \\
    \pointCyan $10^{-a} = \dfrac{1}{10^a}$
    \qq{} \exemple $\dfrac{10^2}{10^1} = 10^2 \times 10^{-1} = 10^{2 - 1} = 10^1$
  \end{importants}
\end{doc}


\numeroQuestion Réaliser les calculs suivants en donnant le résultat sous la forme d'une puissance de dix :
\pasCorrection{\vspace*{-4pt}}
\begin{listePoints}[2]
  \setlength\itemsep{-4pt}
  \item $\num{e2} \times \num{e1}  =$ \texteTrou(0){$10^{2 + 1} = 10^3 = 1000$}
  \item $\num{e4} \times \num{e-3} =$ \texteTrou(0){$10^{4 - 3} = 10^1 = 10$}
  \item $\num{e-2}\times \num{e-3} =$ \texteTrou(0){$10^{-2 -3} = 10^{-5} = 0,00001$}
  \item $\num{e-1}\times \num{e-5}\times \num{e4}=$ \texteTrou(0){$10^{-1 -5 +4} = 10^{-2} = 0,01$}
\end{listePoints}


%%%%
%\newpage
\vspace*{-12pt}
\titreSection{Notation scientifique}

\begin{doc}{La notation scientifique}
  \begin{importants}
  La \important{notation scientifique} d'une quantité se présente de la façon suivante :
  % Textes
  \begin{center}
    \texteEncadre{un chiffre différent de zéro}
    \qq{}
    \texteEncadre{autres chiffres} 
    \qq{}
    \texteEncadre{puissance de dix}
    \texteEncadre{\important{unité}}
  \end{center}
  % Virgule et multiplication
  \vspace*{-30pt} \hspace*{4pt}
  \begin{tikzpicture}  
    \draw [white] (0,0) circle;
    \draw [couleurSec-900, thick] (5.9,-0.25) circle [radius=0.3] node {$\mathbf{,}$};
    \draw [couleurSec-900, thick] (10,0) circle [radius=0.3] node {$\mathbf{\times}$};
  \end{tikzpicture}
\end{importants}

  
  Concrètement, il va falloir compter combien de zéro on doit enlever ou ajouter pour que le nombre soit compris entre 0 et 10, si on en enlève la puissance de 10 est positive, sinon elle est négative.
\end{doc}

\numeroQuestion Écrire les nombres suivants comme le produit d'un nombre compris entre 0 et 10 et d'une puissance de 10 (\exemple \num{600} = \num{6,00e2}) :
\pasCorrection{\vspace*{-4pt}}
\begin{listePoints}[2]
  \setlength\itemsep{-4pt}
  \item \num{100000}  = \texteTrou(0){\num{1e5}}
  \item \num{9000000} = \texteTrou(0){\num{9e6}}
  \item \num{0,0006}  = \texteTrou(0){\num{1e-4}}
  \item \num{0,00705} = \texteTrou(0){\num{7,05e-3}}
\end{listePoints}

\numeroQuestion Écrire les quantités suivantes en notation scientifique :
\begin{listePoints}[2]
  \item \qty{288}{\hour}      = \texteTrou(0){\qty{2,88e2}{\hour}\\}
  \item \qty{1}{\m}           = \texteTrou(0){\qty{1e0}{\m}\\}
  \item \qty{756 864 000}{\s} = \texteTrou(0){\qty{7,56 864 000}{8\s}\\}
  \item \qty{0,01}{\percent}  = \texteTrou(0){\qty{1,0e-2}{\percent}\\}
  \item \qty{8960}{\g/\l}     = \texteTrou(0){\qty{8,96e3}{\g/\l}\\}
  \item \qty{0,436}{\s}       = \texteTrou(0){\qty{4,36e-1}{\s}\\}
\end{listePoints}

\attention Il faut \important{toujours} préciser \important{l'unité} d'une grandeur quand on réalise un calcul !
L'unité indique comment la grandeur a été mesurée.
Sans unités le résultat \important{n'a pas de sens.}
Seuls les rapports entre grandeurs sont sans unités, et on les utilise rarement en physique-chimie en ST2S.
